\section{Constrained paths and their detection estimates}\label{ap:paths}

At one layer of a triangular cell, we seek many integer-affine maps
$\psi(t)=x+Vt$ along which the current polynomial has a small
polynomial residual and an exact integer-polynomial lift. The affine
maps preserve arithmetic progressions; the polynomial identity will
preserve the determining equations when a new patch is returned to
the parent. We construct their probability law in this section.

There are two comparison estimates to prove for that law. With the
center fixed, tests chosen separately on the sampled boxes must be
detectable on the parent. After averaging the center, the one-site
distribution must approximate the parent distribution. Below the
current degree, the first comparison must be measured relative to
the small chart: charging its reciprocal mass to the output
complexity would charge the current width to its own relative loss.

The coefficient density will be a product law on an abstract
coefficient torus. The actual ensemble is the integer spatial law
weighted by that density. These are different probability spaces.
The construction below defines both; high rank will supply the
comparison between them. The general estimates for coefficient
modes and polynomial images are proved in
Appendices~\ref{sampling:constrained-section} and~\ref{cube:section}.

Throughout this section, $j$ is the current layer, $d=d_j$,
$W\subseteq\R^d$ is rational, and
\[
 U=\prod_{a=1}^n\{0,\ldots,H_a-1\},\qquad
 C: \R^n\longrightarrow W,\qquad \deg C\le j.
\]
Lower layers have already been substituted.  Rank means the ordinary
rank condition of Definition~\ref{sampling:rank}, with the single tag
$j$: no integer row of norm at most $R_j$, nonzero on $W$, takes the
degree-$j$ symbol of $C$ to a slow-plus-rational polynomial at bound
$R_j$.  Warm, cold-pre, and late choices have the meanings fixed in
Definition~\ref{ap:budgets}.  All degrees and cube orders below belong to a
fixed finite range.  Structural numbers of columns are polynomial in
dimensions and are fixed before numerical precisions.

\subsection{Normalized lift charts}

Put $\Lambda=W\cap\Z^d$ and
$\Lambda'=\operatorname{proj}_{W^\perp}\Z^d$.
Choose the basis $v_i$, $1\le i\le\dim W^\perp$, of the projected
lattice $\Lambda'$ supplied by Lemma~\ref{sampling:lattice-chart},
and set
\begin{equation}\label{ap:path:mass}
 K_i=\max(1,\lceil |v_i|^{-1}\rceil),\qquad
 k_i=\max(1,\lceil K_i v\rceil),
 \qquad
 m_j=v^{\dim W}\frac{\prod_i k_i}{\covol(\Lambda)}.
\end{equation}
Here $v$ is a sufficiently small warm fraction of $w_j$.  All lift
coordinates use this same physical width.  More precisely, $v/w_j$ has
logarithmic cost polynomial in dimensions and $\log\gamma^{-1}$.
Principal amplitudes in normalized lift units, and the proportionality
constants in inactive side lengths, can have logarithmic costs depending
only on dimensions.  They are chosen small enough that all true and
ideal lifts lie well inside a buffered chart of radius comparable to
$v$, itself inside the required gap.  These charts are narrow enough
that every Boolean difference of an order used below has components
of magnitude less than $1/2$.

For a center $c\in W/\Lambda$, choose a representative in $W$.  On a
uniqueness chart write
\[
 C(u)-c=y(u)+\beta(u),\qquad \beta(u)\in\Z^d.
\]
The sheet coordinates satisfy
\[
 \operatorname{proj}_{W^\perp}y=\sum_i v_i z_i
       =-\operatorname{proj}_{W^\perp}\beta.
\]
Choose integer lifts of the basis vectors $v_i$, and let $l(z)$
be their integer linear combination. Then $\beta+l(z)\in\Lambda$.
Choose an integer basis of $\Lambda$ and let $w$ be the coordinates
of $\beta+l(z)$ in that basis. Primitivity of $\Lambda$ shows that
$\beta\mapsto(z,w)$ is an integer unimodular change of coordinates.
Here $w$ is a deck coordinate, not a width. A residue character of
$(z,w)$ is therefore a residue character of $\beta$ at the same
modulus; the real size of the basis coefficients does not enter
this assertion.

\begin{lemma}[Chart mass and normalized coordinates]\label{ap:path:chart}
The quantities $m_j$ and $m_j^{-1}$ have cold-pre bounds independent of
the height of $W$.  On a buffered chart the coordinates
$y_W/v$ and $z_i/k_i$ have dimension-controlled norms in terms of $y/v$.
The row defining $z_i$ on the ambient lift has norm at most
$\exp(\operatorname{poly}(d))K_i$, is integral on $\Z^d$, and
annihilates $W$.

A bounded warm enlargement of the chart has Haar measure at most a
warm multiple of $m_j$.  A ball of a warm fraction of radius $v$,
centered at a feasible torus point and with the prescribed room, has
Haar measure at least a warm inverse multiple of $m_j$.  These relative
bounds do not depend on any width when the radius ratios use only
dimensions and $\gamma$.
\end{lemma}

\begin{proof}
Lemma~\ref{sampling:lattice-chart} gives
\[
 \covol(\Lambda)\asymp_{\exp(\operatorname{poly}(d))}\prod_iK_i
\]
and dimension-controlled maps between the projected Euclidean
coordinates and $z_i/K_i$.  Since
$k_i\asymp\max(1,K_i v)$, multiplying the latter coordinates by
$K_i v/k_i\le1$ proves the asserted upper norm bound.  The dual
projected lattice is the integral annihilator of $W$, giving the
integral rows and their stated norms.  If $q=\dim W^\perp$, then
\[
 \exp(-\operatorname{poly}(d))v^d
       \le m_j\le \exp(\operatorname{poly}(d))v^{\dim W}.
\]
In particular its logarithmic costs involve $\log v^{-1}$ but not
subspace height.

Haar measure in these coordinates is Lebesgue measure on $W$ times
counting measure on the sheets, divided by $\covol(\Lambda)$.
Counting sheet coordinates in a dimension-controlled parallelepiped
and integrating on $W$ gives the upper bound.  Around a feasible
sheet, a smaller such parallelepiped contains at least a
dimension-controlled multiple of $\prod_i\max(1,\tau k_i)$ sheet
points when its radius is $\tau v$; this is at least
$\tau^q\prod_i k_i$ up to dimension factors.  Integration on $W$
supplies the other $\tau^{\dim W}$ factor.  This proves the lower
bound, including arbitrary fixed real translations of the chart.
Feasibility is necessary: no assertion is made that an arbitrary
translated small box meets the torus.
\end{proof}

\subsection{The coefficient law and the actual path ensemble}

First consider the model $W=\R^d$, where there are no discrete
sheet coordinates. We want the spatial directions and the small-lift
directions to be recoverable separately. The reference lift law
therefore uses disjoint blocks of $j$ parameters with a common
side $L$, to be selected later. On a lift axis $a$, its principal
part has the form
\[
 \frac{y_a^{\rm ideal}(t)}v=\eta_a+
       \sum_B c_{aB}\prod_{\ell\in B}\frac{t_\ell}{L},
\]
with the normalized $c_{aB}$ bounded away from zero. The constant
$\eta_a$ supplies the local translation. The blocks for different
axes are disjoint, and a separate kernel block controls the spatial
image. Fixing all but one variable in each selected product at
normalized values bounded away from zero gives independent
directions in this ideal model. Repeating the products
on further disjoint blocks also smooths their one-site and cube
evaluation laws.

The reason to recover lift directions rather than raw integer
coefficients comes from the lift identity we require. We will
arrange that $C(\psi(t))-c=y(t)+\beta(t)$ is a polynomial identity
of degree at most $j$. Its top homogeneous parts then give
\[
 \beta_{[j]}(t)-C_{[j]}(\psi_{[1]}(t))=-y_{[j]}(t).
\]
Here brackets denote ordinary homogeneous degree. Thus the products control the deviation of the integer variable
from its prescribed polynomial value without any bound on the
coefficients of $C$. Small remaining lift coefficients and small
nonkernel spatial displacements perturb this ideal recovery map;
they do not give exact independence to the full map.

For a rational $W$, the projected lift coordinates are integer
sheets. Their available ranges determine whether the preceding
products can use long parameters or must instead use coefficient
one and shorter parameters. We now prescribe those ranges and
the additional columns needed for residue tests and exact sheet
translations. All independence in this prescription belongs to
the reference coefficient law and the separate spatial law.
The actual path law couples them by the tilt below.

After the pre requirements are fixed, choose a perturbation
parameter $0<\sigma<1$ and a range parameter $A\ge2$, in that
order. Rank thresholds and parent sides will be chosen later.
Choose a scale gap with no $k_i^{1/j}$ between $L$ and $AL$.
For example, among more than $\dim W^\perp$ consecutive disjoint
multiplicative intervals of ratio $A$, at least one contains none
of these numbers. Hence this choice is possible with
$A\le L\le A^{O(1+d)}$.  With a fixed sufficiently large
degree-dependent $C_*$, call a sheet axis
\[
 \begin{array}{ll}
 \text{inactive}& k_i\le L^j,\\
 \text{moderate}& (AL)^j\le k_i\le L^{C_*},\\
 \text{enormous}& k_i>L^{C_*}.
 \end{array}
\]
There are only polynomially many numbers to avoid, so $\log L$ can
be chosen with the permitted late bound without using their magnitudes.
Rounding is performed with fixed buffers at the cutoffs.
An active sheet axis means a moderate or enormous axis.

The parameter box is
$\mathcal T=\prod_\ell\{0,\ldots,T_\ell-1\}$.
We prescribe its columns, the independent reference law of the
lift coefficients, and the supports of the separate spatial law
as follows. The actual coefficients of $C(x+Vt)$ are determined
by $(x,V)$; the tilt defined below combines these two laws.
\begin{enumerate}
\item To keep the spatial image nondegenerate independently of the
principal lift directions, spatial kernel columns have side $L$.
In normalized spatial coordinates their matrices lie in buffered neighborhoods of fixed
invertible matrices, at a common warm spatial scale.  The inverse
norms are warm-bounded throughout the supports.
\item Continuous and active lift axes have independent dedicated
blocks of $j$ columns, all of side $L$.  Their principal monomials
are squarefree products on disjoint blocks.  In normalized units
$y_W/v$ and $z_i/k_i$, the principal amplitudes have buffered
supports bounded away from zero at the chosen dimension scale.
Equivalently, a principal raw coefficient has size comparable to
$v/L^j$ on a continuous axis and $k_i/L^j$ on an active sheet
axis, with the corresponding smooth continuous or integer law.
\item On an inactive axis, integral coefficients cannot be made
arbitrarily small. We therefore keep the product coefficient
exactly one and shorten the dedicated sides to quantities comparable
to $k_i^{1/j}$ at the prescribed dimension factors.  If these sides would be below a specified
warm minimum, omit the axis: set \emph{every} coefficient, including
the constant, equal to zero, and supply no columns for it.  Hence
omitted $k_i$ are warm-bounded.  Otherwise supply one additional
nonkernel column $t_{i,0}$, with side comparable to $k_i$.
Its coefficient is one so that translating it translates the
sheet output by exactly the same amount, as needed for the
later affine sections. Set
\begin{equation}\label{ap:path:inactive-polynomial}
 z_i(t)=z_i(0)+t_{i,0}
                 +\sum_{\text{dedicated blocks }B}\prod_{\ell\in B}t_\ell.
\end{equation}
There are no other terms on this axis.  The common proportionality
constants are small enough for chart containment.
\item To test congruence relations without sharing the principal
smoothing variables, add a structural polynomial number of free
nonkernel columns, of side $L$, disjoint from the preceding blocks.
Their independent leading coefficient data supply the
residue-test vectors in the reference comparison law.  On moderate
axes, put independent small perturbations on all pure-free
nonconstant monomials, with integer coefficient widths comparable to
$k_i\sigma/T^\nu$.  Other nondesignated moderate coefficients remain
zero. The scale gap makes the selected integer ranges broad
enough for later individual-modulus tests. On continuous and
enormous axes, each remaining nonconstant
monomial $t^\nu$ has an independent small coefficient whose width
in the normalized output coordinates is comparable to
$\sigma/T^\nu$. The continuous coefficients have smooth laws.
The enormous-axis coefficients are integers, with the smooth
discrete laws specified below.
\end{enumerate}
Here and below fixed structural factors may multiply these widths.
Every constant lift coefficient has an independent law of fixed
warm support in its normalized units. When its integer range is
too short for that law it is fixed at zero; its corresponding
$k_i$ is then warm-bounded. Counts of kernel and product blocks
are fixed polynomials in the structural dimensions, large enough
for the fixed cube orders and smoothing orders used below, with
spare blocks for smoothing. Numerical precisions never increase
these counts. Every nonomitted side exceeds the specified warm
minimum, which may include higher-layer parent and separation
requirements. Every side is $O(L^j)$.

We specify the coefficient laws as well as their supports.
Fix a nonnegative, nonzero smooth bump $\varphi$ supported in
$(-1,1)$ and positive on a neighborhood of zero.
An integer coefficient with center $a$ and scale $S$ has law
\[
 \Pr(\alpha=k)=
 \frac{\varphi((k-a)/S)}
      {\sum_{r\in\Z}\varphi((r-a)/S)}.
\]
The center and scale place this support strictly inside the
prescribed interval and ensure that its positive region contains
an integer, so the denominator is positive. Coefficients fixed at zero or one have their
stated point-mass laws. Continuous coefficients are their centers
plus their prescribed widths times independent variables with
fixed smooth buffered densities. All these choices are independent
across coefficient positions and output axes.

For a monomial $t^\nu$, let $f_\nu$ be the resulting density of its
coefficient in the orthogonal $W$ coordinates, and let
$p_{i,\nu}$ be the probability mass of its coefficient on sheet
axis $i$. On the coefficient torus $W/\Lambda$, the local Haar
measure is $du_W\,d\#z/\covol(\Lambda)$. Consequently
\[
 \covol(\Lambda)\,f_\nu(u_W)\prod_i p_{i,\nu}(z_i)
\]
is the density of this coefficient law relative to Haar.
The selected supports lie in a buffered coefficient chart.
Their product over $|\nu|\le j$ defines a probability density
$\mathcal D_c$ on
$\prod_{|\nu|\le j}W/\Lambda$; the center $c$ shifts only the
constant coefficient. In particular $\int\mathcal D_c=1$.

We also choose a law $\rho(x,V)$ on integer spatial coefficients.
The root $x$ is nearly uniform on a buffered interior box, with
a smooth transition near its boundary. In each spatial kernel
block the normalized matrix lies in a product neighborhood of a
fixed invertible matrix at the selected warm spatial scale.
Its integer entries are sampled independently by smooth buffered
bumps in their normalized intervals, with the same normalization
as the integer coefficient law above. Every nonkernel entry is
independent and has a smooth profile
supported in
\[
             |V_{a\ell}|
             \le \frac{c_0\sigma H_a}{mT_\ell},
\]
where $m$ is the number of columns and $c_0>0$ is a fixed small
constant. These widths become long integer intervals when the
parent sides are chosen. Their total normalized displacement
is $O(\sigma)$. All supports have enough room that every path
lies in $U$.

The actual law of the affine maps is
\begin{equation}\label{ap:path:tilt}
 d\mathbb P_c(\psi)
  =\frac{\rho(x,V)\mathcal D_c(\operatorname{coeff}C(x+Vt))}
  {Z_c},\qquad \psi(t)=x+Vt.
\end{equation}
The normalizer in this formula is
$Z_c=\E_\rho\mathcal D_c(\operatorname{coeff}C(x+Vt))$.
The identity $\int\mathcal D_c=1$ concerns abstract Haar coefficients;
the conclusion $Z_c=1+o(1)$ will use high rank of the actual
polynomial coefficients.

There is a second distinction when residue information is needed.
On the finite cover $W/M\Lambda$, Haar measure above a coefficient
in $W/\Lambda$ is uniform on its $M^{\dim W}$ deck choices.
Thus the deck coefficients modulo $M$ are independent uniform
variables conditional on the small-lift coefficient data under
this comparison law. The actual integer spatial law will only
approximate the resulting joint residue law at each prescribed
individual modulus. It need not have this exact conditional
independence after real coefficient values are fixed.

By contrast, the lift identity on a path is exact. On the support
of $\mathcal D_c$, subtracting the small real coefficient lifts
from $C(x+Vt)-c$ leaves integer monomial coefficients. Their sum
is a polynomial $\beta(t)$, and the small coefficients sum to
$y(t)$. Chart containment identifies these polynomials with the
site lifts throughout the path.

We will use high rank through the following Fourier estimate.
Consider coefficient characters on a prescribed integer cover,
represented by rational ambient rows $\lambda_\nu$ whose
denominators divide its modulus. For fixed parameter sites
$t_\omega$, evaluation sends a
$W$-valued polynomial coefficient array $(a_\nu)_\nu$ to
$(\sum_\nu a_\nu t_\omega^\nu)_\omega$. A coefficient mode
$e(\sum_\nu\lambda_\nu a_\nu)$ \emph{factors through these sites}
if its real linear functional is a linear combination of the
evaluation functionals; equivalently, it vanishes on their
common kernel. Here $e(z)=\exp(2\pi i z)$.

For one site, or sites
$t_\omega=t_0+\sum_{i=1}^q\omega_it_i$ forming an affine Boolean
cube of fixed order with independent differences, Lemma~\ref{sampling:analytic-removal}
removes every mode which does not factor through the sites:
\[
 \left|\E_\rho
 e\!\left(\sum_\nu\lambda_\nu
          \operatorname{coeff}_\nu(C(x+Vt))\right)
       \prod_\omega F_\omega(x+Vt_\omega)\right|
 \le\varepsilon
 \qquad(|F_\omega|\le1).
\]
The rank and parent-side thresholds depend polynomially in
logarithmic scale on the prescribed Fourier frequencies and
coefficient mass, cover modulus, spatial smoothness, inverse
error, $\log L$, and $\log\sigma^{-1}$. They do not depend on the
height of $W$. Spatial residue conditions may be included.
With no site insertions, every nontrivial coefficient character
is removed. The estimate applies to the Fourier expansion of
$\mathcal D_c$ before dividing by $Z_c$.

Factoring through real evaluation need not give a character on
the original evaluation torus. Clearing the denominators of a
rational right inverse gives a finite cover on which it does.
The cube comparison below retains this cover. For a single site
the constant coefficient already gives an integral right inverse.
These observations specify the estimate used in the following
lemma; its proof remains in the appendix.

\begin{lemma}[The normalized one-layer ensemble]\label{ap:path:ensemble}
For the law just defined, sufficiently large permitted rank and parent
side cutoffs give the following conclusions, uniformly in the center.
\begin{enumerate}
\item $Z_c=1+o(1)$, and on every retained path
\begin{equation}\label{ap:path:exact-lift}
 C(\psi(t))-c=y(\psi(t))+\beta(\psi(t))
\end{equation}
is a polynomial identity.  The small lift is fully slow at width $v$;
$\beta\circ\psi$ has integral monomial coefficients.
\item The non-site-mode removal of
Lemma~\ref{sampling:analytic-removal}, and the Haar and individual-cover
conclusions of Corollaries~\ref{sampling:normalization}
and~\ref{sampling:finite-residues}, hold with cold-pre requests and
late errors.  In particular, on each prescribed individual modulus,
the relevant coefficient residues have the product chart/deck law,
independently of the unweighted spatial residues, to arbitrarily small
prescribed total variation.  On moderate axes the pure-free
nonconstant coefficients are independently uniform to this accuracy
once the range parameter is chosen sufficiently large.
\item Noninjective paths can be discarded at any prescribed cold-pre
probability.  At fixed center the averaged one-site law is supported
on the buffered chart and has density at most $B/m_j$ relative to
uniform measure on $U$, apart from arbitrarily small excess mass,
where $B$ is warm-bounded.
\item Given a warm spatial error $\epsilon>0$, the root law can be
chosen so that Haar mixing of centers gives the one-site upper bound
$(1+\epsilon)$ times the uniform law, and the corresponding lower
bound $(1-\epsilon)$ away from a warm-thickness boundary strip,
each with arbitrarily small cold-pre additive error against every
nonnegative function bounded by one.  The spatial scale and strip
thickness do not depend on $Q_j$.
\end{enumerate}
\end{lemma}

\begin{proof}
The independent coefficient density has Haar integral one.  On its
support subtracting the specified small coefficient lifts leaves
integral coefficients.  Summing the resulting monomials proves
Equation~\eqref{ap:path:exact-lift}; uniqueness of the chart identifies
this polynomial lift with the site lift throughout the path.

We check the quantitative hypotheses for Fourier removal rather than
deduce $Z_c=1$ from the abstract normalization.  In a coefficient
chart the Haar density has the factor
$\covol(\Lambda)/(v^{\dim W}\prod_i k_i)=m_j^{-1}$.
The remaining factors are normalized continuous densities and
normalized discrete bumps.  On axes with $k_i\le L^{C_*}$, singleton
interpolation costs only $O(\log L)$ in addition to cold-pre costs.
This includes omitted axes.  On enormous axes all coefficient ranges
contain many grid points.  Since $T^\nu\le L^{j^2}$ for degrees at
most $j$, their normalized smooth interpolation has logarithmic costs
polynomial in the cold-pre budget, $\log\sigma^{-1}$, and $\log L$,
when $C_*$ is sufficiently large.  Lemma~\ref{ap:path:chart} controls
all normalized coordinate maps independently of orientation and
height.  Buffered ambient chart extensions can therefore be
periodized and Fourier approximated with just these costs.

On a modulus cover multiply by the residue character of the chart
index $\beta$.  Rows defining $(z,w)$ reduce modulo that modulus,
so no real deck-basis norm is needed.  The same ambient extension
works on the cover with its individual logarithmic modulus cost.
The displacement proof of Lemma~\ref{sampling:analytic-removal}
only changes by replacing its parameter bound $L$ with $O(L^j)$.
Its translating boxes and the parent cutoff still have permitted
late logarithmic costs.  Thus it removes each non-site mode with
arbitrary bounded site insertions.  With no insertions it removes
every nontrivial coefficient mode; the trivial mode is the Haar
integral one.  This proves $Z_c=1+o(1)$.  Translating $c$ changes
Fourier phases but none of the bounds, proving uniformity.

For each prescribed residue vector apply this argument on its own
cover and sum per-atom errors.  The logarithm of the atom count is
polynomial in dimensions and the individual modulus log.  This is
the total-variation argument of
Corollary~\ref{sampling:finite-residues}, with the ambient extension
just verified.  A moderate pure-free coefficient of top degree has
width at least a structural factor times $A^j\sigma$; coefficients
of smaller degree have no smaller range.  Choosing $A$ after the
individual-modulus requests makes every such smooth integer law
uniform modulo that modulus to the required accuracy.  Parent sides
similarly make spatial free entries broad.  The deck laws are
uniform independently of the sheet data.  This proves the full
joint assertion, not merely a claim about designated top monomials.

The independent constant coefficient has Haar density at most
$B/m_j$ for a warm $B$.  The surviving one-site weight under
non-site removal is its convolution with the probability law of
the other coefficient contributions, so has the same cap.  The
independent spatial root has a warm cap conditional on all slopes.
Removal with arbitrary site indicators consequently bounds the
excess mass above a fixed warm multiple of $m_j^{-1}$ by the
requested error.  All actual supports are buffered.

When centers are mixed, a mode factoring through one site has its
site row as its constant coefficient row.  Center integration kills
every nontrivial such row; removal kills the nonfactoring modes.
Only the unweighted spatial law remains.  Choose its root nearly
uniform on an interior box and make all spatial displacements a
smaller warm fraction of the sides.  Its translate has the asserted
upper density everywhere and lower density off the indicated strip.
Only the Fourier removal error, not that strip, is requested at
cold precision.

Finally, for a fixed nonzero difference $a=t-t'$, the equation
$Va=0$ determines at most one value of an entry in a column for
which $a_\ell\ne0$.  Smooth broad spatial laws, a union bound over
at most $O(L^{2jm})$ pairs, and the permitted late supremum bound
of the tilt make collisions negligible by increasing parent sides.
This completes all the claims in the permitted choice order.
\end{proof}

\subsection{Comparing sampled and ambient tests}

\subsubsection{Niltests}\label{ap:niltest-conventions}
A rational filtered nilmanifold of degree $q$ consists of a connected,
simply connected nilpotent Lie group $G$, a cocompact lattice $\Gamma$,
and connected rational subgroups
\[
 G=G_1\supseteq\cdots\supseteq G_{q+1}=\{1\},
 \qquad [G_i,G_j]\subseteq G_{i+j}.
\]
Rationality is measured in fixed exponential coordinates on the Lie
algebra $\mathfrak g=\log G$, with rational bases for
$\mathfrak g_i=\log G_i$. A polynomial map $g:\Z^n\to G$ is
adapted when $\log g$ modulo $\mathfrak g_{i+1}$ has degree at
most $i$, using weighted degree when variable weights are assigned.
A \emph{niltest} is the function $u\mapsto F(g(u)\Gamma)$,
where $F$ is bounded and Lipschitz on $G/\Gamma$. The metric is
the quotient of the right-invariant Riemannian metric making the
chosen coordinate basis orthonormal at the identity.

Complexity at most $p$ means that $\dim G$, the logarithmic heights
of the structural constants and filtered bases, and
$\log(2+\|F\|_\infty+\Lip F)$ are at most $p$, with
\[
 l\Z^{\dim G}\subseteq\log\Gamma\subseteq l^{-1}\Z^{\dim G}
 \quad\text{for some integer }1\le l\le e^p.
\]
A metric change includes the logarithm of its comparison constant in
this budget. No coefficient bound is imposed on the orbit $g$.
On $\Z/N\Z$, evaluation uses representatives $0,\ldots,N-1$;
degree-zero tests are constants, and dimension zero is allowed.
Definition~\ref{prelim:niltest} repeats these conventions where the
rational calculus is developed. The equivalence with adapted
Taylor/binomial polynomial sequences is proved there.


\subsubsection{Cube normalization}\label{ap:cube-conventions}
For a finite integer box $Q$, let $\operatorname{Cub}_j(Q)$ be the
set of tuples $(x,h_1,\ldots,h_j)$ for which every vertex
$x+\omega\cdot h$, $\omega\in\{0,1\}^j$, belongs to $Q$.
Writing $\mathcal C z=\overline z$, we use
\[
 \|f\|_{U^j(Q)}^{2^j}
 =\frac{1}{|\operatorname{Cub}_j(Q)|}
   \sum_{(x,h)\in\operatorname{Cub}_j(Q)}
   \prod_{\omega\in\{0,1\}^j}
       \mathcal C^{|\omega|}f(x+\omega\cdot h).
\]
On a finite abelian group, all parameters instead have independent
uniform averages. In particular, the constant function $1$ has norm
one under both conventions. Formula~\eqref{prelim:box-gowers}
and the identities following it give the exact comparison with
zero extension to a product of cyclic groups and the cube-count
weights required when taking derivatives.




A \emph{normalized site twist} is a unit-bounded function localized
to a buffered chart which is a product, or a bounded-complexity
combination of products, of a smooth function of $(u/H,y/v)$ and a
bounded-modulus function of $(u,\beta)$.  Smoothness is measured in
these normalized variables.  Such a twist need not have warm
Lipschitz cost as a function on the unlocalized ambient torus.

\begin{lemma}[Cube comparison in chart units]\label{ap:path:comparison}
Fix slice costs, a fixed cube order $q$, and a requested accuracy.
For a fixed permitted parameter slice $B$, the sampled functional
\[
 (P_\omega)_\omega\longmapsto
 \E_\psi\E_{\mathbf t\in\operatorname{Cub}_q(B)}
            \prod_{\omega\in\{0,1\}^q}P_\omega(\psi(t_\omega))
\]
is, uniformly for $|P_\omega|\le1$, approximated to the requested
accuracy by
\[
 \sum_\alpha c_\alpha
 \E_{\mathbf u\in\operatorname{Cub}_q(U)}
       \prod_{\omega\in\{0,1\}^q}
                P_\omega(u_\omega)J_{\alpha,\omega}(u_\omega).
\]
The site twists $J_{\alpha,\omega}$, coefficient mass
$\sum_\alpha|c_\alpha|$, and number of terms have cold-pre bounds.  If $q\le j$ and
the request is warm, its ambient output can instead be written
\begin{equation}\label{ap:path:sparse-comparison}
 \sum_\alpha c_\alpha\,
 \E_{\text{ambient cubes}}
       \prod_{\omega\in\{0,1\}^q}
       P_\omega(u_\omega)\frac{J_{\alpha,\omega}(u_\omega)}{m_j},
\end{equation}
with warm total coefficient mass, warm normalized site twists
$J_{\alpha,\omega}$, and an arbitrarily small prescribed warm
additive error, uniformly for arbitrary $|P_\omega|\le1$.
Finite sums suffice.  Only the error transfer and removal, not the
output twists or coefficient mass, use cold or late costs.
\end{lemma}


The proof appears after the general cube comparison in
Appendix~\ref{cube:section}. Its factors depend only on individual
ambient sites. This is the feature that lets a large sampled cube
average yield an ambient niltest correlation.

Choose a nonnegative smooth cutoff $\chi$ supported in a buffered
uniqueness chart and equal to one on all path and twist supports.
For fixed center $c$, normalized chart measure on $U$ is the
probability measure
\[
 \mu_c(\{u\})=
 \frac{\chi(C(u)-c)}{\sum_{v\in U}\chi(C(v)-c)}.
\]
The denominator is positive for the sufficiently large rank and
side choices below. This measure is used only for the warm
detection assertion; the approximation seminorm is defined from
the actual path law.

Fix a degree $r$, permitted slice costs, and local test complexities.
For a site input $g:U\to\C$, put
\[
 X_r(g)=\E_{\psi\sim\mathbb P_c}
        \sup_{B,F}\left|\E_{t\in B}g(\psi(t))F(t)\right|,
\]
where $B$ ranges over the permitted slices of $\mathcal T$ and $F$
over unit-bounded local niltests of degree $r$ and the prescribed
complexity. Each path may choose its own slice and niltest. An
ambient atom is a product $J(u)N(u)$, where $J$ is a normalized
site twist and $N$ is an ordinary ambient degree-$r$ niltest;
both factors may be taken unit-bounded.

\begin{theorem}[Pathwise detection and approximation]
\label{ap:path:detection}
Prescribe finitely many testing degrees in the fixed allowed range,
local test complexities, slice costs, input caps, and accuracies.
The later parameters of the one-layer construction can be chosen
once so that the following assertions hold for this same ensemble,
uniformly in the center and the site input.

For cold-pre requests, if a cold-pre-bounded input satisfies
$X_r(g)\ge a>0$, with $\log(1/a)$ bounded by cold-pre data,
then $g$ correlates in uniform parent measure with an ambient
atom of cold-pre complexity, and the logarithm of the reciprocal
correlation has a cold-pre bound. Every $|g|\le1$ can also
be approximated in $X_r$, to the prescribed cold-pre accuracy,
by a linear combination of such atoms with cold-pre coefficient
mass and number of terms.

If $r<j$ and the requests are warm, every warm-bounded input
with $X_r(g)\ge a>0$ and a warm bound for $\log(1/a)$
correlates in $\mu_c$ with an ambient atom of warm complexity,
with a warm bound for the logarithm of the reciprocal correlation.
The atoms, coefficient mass, and number of terms in the
$X_r$ approximation also have warm bounds. These
bounds may use the omission threshold but do not use $Q_j$ or
the height of $W$.

For a specified larger input cap, include its logarithm in the
request budget and rescale. Increasing the requested accuracy,
test, or path budgets changes the corresponding warm or cold-pre
output logarithms polynomially. All path probabilities and
testing suprema refer to the law before productive paths are
selected.
\end{theorem}

\begin{proof}[Proof of the cold clause]
The ordinary argument of Lemma~\ref{cube:pathwise-detection}
applies to the ensemble above by Lemma~\ref{ap:path:comparison}.  Indeed, mesh endpoints and enumerate strides to reduce
the separately chosen slices to finitely many possibilities.  A
short side has a bounded number of points and can be enumerated
directly.  A positive fraction of detecting paths therefore has
a large $U^{r+1}$ norm on one common slice.  The niltests on these
paths need not be the same.

Lemma~\ref{ap:path:comparison} transfers their averaged cube
power to a mixed ambient cube correlation.  When all physical
chart costs are included in the cold preliminary budget, its
site factors have controlled caps.  Gowers--Cauchy--Schwarz and
Theorem~\ref{prelim:inverse}, as in the cited detection proof,
give an ambient degree-$r$ niltest times a permitted site twist
correlating with the input.  Degree zero uses the two independent
sites of a one-cube and gives a twist alone.

To pass from detection to approximation, let $m_U$ be the full
one-site density relative to uniform measure on $U$.  If every
permitted slice has at least fraction $\delta_{\rm sl}$ of the
full parameter box, then
\[
              X_r(\phi)\le
       \delta_{\rm sl}^{-1}\mathbb E_U m_U|\phi|.
\]
Lemma~\ref{ap:path:ensemble} bounds the excess mass above a cap
of cold preliminary size.  Doubling the cap bounds the mass of
the resulting tail by twice that excess.  Select this cap first,
then request the tail accuracy needed by Lemma~\ref{cube:separation}.  That
lemma produces the finite approximation, with controlled atom
count and coefficient mass.  All requests precede the late
rank and side choices.  This proves the cold clause for the
same ensemble as the warm clause. Section~\ref{ap:relative-detection}
proves its bounds independent of the current chart mass.
\end{proof}

\subsection{The geometry needed in lifting}

There is a deliberate distinction from the packaged geometry in
Property~\ref{sampling:property-geometry}.  Mixed-center spatial
comparison here has a warm-thickness boundary strip; it is not
claimed with a strip as thin as every cold detection accuracy.
This distinction does not change the path deductions used below.
In the proof of Proposition~\ref{lifting:candidate-major}, local
correlations are generated at fixed centers by freezing pre-short
axes and using a chart-supported phase input.  The rest of that
argument uses the exact fully slow/integral identities,
fixed-center testing detection, translation groups, and the
ambient rank estimate.  Lemma~\ref{ap:path:ensemble} and the
cold clause of Theorem~\ref{ap:path:detection} provide precisely
these data; physical-coordinate or cover twists then have
cold-pre logarithmic Lipschitz and modulus costs.  Changing the
upper parameter side bound to $O(L^j)$ changes only late
cutoffs, and omitting axes removes variables from these arguments.
The same inspection applies to the symbol and modeling steps in
Proposition~\ref{lifting:candidate-globalization}.  For passive
partner selection below, it is the warm \emph{approximation}
clause of Theorem~\ref{ap:path:detection}, rather than a
whole-box detector of warm mass, that will be used.


\subsection{The stopped tree of a triangular cell}
\label{ap:path:stopped-tree}

Fix the current triangular cell, its value spaces $W_j$, and one
collection of numerical requests for all uses of the one-layer
ensemble.  We now define the probability measure in which preparation
will test rank failures.  The layers are sampled in the order
$1,\ldots,D$, although their numerical budgets are chosen in the
opposite order.

At entry to layer $j$, a history consists of the earlier affine path
maps, their centers and their integral lift polynomials.  Substituting
these maps and lifts into $C_j$ gives an ordinary polynomial on the
current parent box, still taking values in $W_j$.  The same exact
substitutions have been made in every higher determining polynomial.
We use $C_j$ for this pulled-back polynomial.

First test the rank of $C_j$ at the current parent scale.  If the test
fails, stop that history.  Otherwise choose its center with Haar
probability conditioned on
\begin{equation}\label{ap:scalar:feasible-center}
 \mathcal C_j=
 \left\{c\in W_j/(W_j\cap\mathbb Z^{d_j}):
       \|c+l_j-b_{j,0}\|_\infty<(1+\gamma/2)w_j
       \text{ for some }b_{j,0}\in\mathbb Z^{d_j}\right\}.
\end{equation}
The condition is independent of the chosen representative of $c$:
changing it by an element of $W_j\cap\mathbb Z^{d_j}$ changes
$b_{j,0}$ by the same integer vector.  Once a representative is fixed,
$b_{j,0}$ is unique.  Sample $\psi$ from the fixed-center ensemble
$\mathbb P_c$.  Its polynomial lift identity is
\[
 C_j(\psi(t))-c=y_j(t)+\beta_j(t),
                    \qquad b_j(t)=\beta_j(t)+b_{j,0}.
\]
The second equation specifies the raw determining slot used in all
higher polynomials.  It gives
\[
 b_j(t)-C_j(\psi(t))-l_j
                   =b_{j,0}-c-l_j-y_j(t).
\]
The lift width is chosen sufficiently small inside the gap that the
enlarged weight-$j$ constraint holds on the whole retained path.
Substitute this affine map and this exact integer polynomial into the
remaining determining blocks and proceed to layer $j+1$.

For applications that discard additional sampled paths, fix the
discard rules as part of the current tree being tested.  At each
history a rule specifies a subset of the finite support of
$\mathbb P_c$; outcomes outside that subset are killed before
proceeding.  In the forward argument these subsets impose injectivity
and the modular conditions defined there, and their lost probability
must satisfy the bounds prescribed in the forward-mass proof.
They delete probability mass; they do not renormalize
the distribution on the remaining histories.  A mask of the input
function, such as a site-dependent flag, is likewise distinct from
conditioning the path law.

Thus the histories arriving at each layer carry a subprobability
measure.  The rank-stop probability at $j$ means the mass of histories
that reach that test and fail it; histories already stopped contribute
zero.  All conditional probabilities used in backward transfer are
ordinary conditional probabilities within this one tree.  After a
rank cut the cell and its polynomials change, so preparation builds
this tree again and repeats the tests.  It does not assume that an
earlier successful rank test survives the cut.

Write $e_j$ for the Haar measure of $\mathcal C_j$.  It depends on
$W_j,l_j,w_j,\gamma$ and not on the sampled prefix.  The strict
certificate ensures that $\mathcal C_j$ is nonempty with positive
measure.  The quantitative lower bound for $e_j$, and the comparison
of strict and enlarged masses down this tree, will be proved with
the forward-mass argument.
